% Diagrama 3. Flujo de la supervisión durante la impresión, en TikZ.
% Equivale a D3_flujo_supervision.mmd. (P) = propuesto, no fijado en la tesis.
% Compilar: pdflatex D3_flujo_supervision_tikz.tex
\documentclass[border=4pt]{standalone}
\usepackage[T1]{fontenc}
\usepackage[utf8]{inputenc}
\usepackage[spanish,es-noshorthands]{babel}
\usepackage{times}
\usepackage{amssymb}
\usepackage{tikz}
% Estilos comunes de los diagramas de validación (TikZ).
% Blanco y negro, Times a 8 pt, sin colores, como la Figura 22 de la tesis.
\usetikzlibrary{positioning,calc,shapes.geometric}
\tikzset{
  base/.style   = {font=\fontsize{8}{9.6}\selectfont},
  proc/.style   = {draw, fill=white, align=center, text width=3.0cm,
                   minimum height=0.9cm, inner sep=3pt,
                   execute at begin node={\hyphenpenalty=10000\relax}},
  propp/.style  = {proc, dashed},
  negra/.style  = {proc, fill=black!12},
  falla/.style  = {proc, fill=black!7},
  dec/.style    = {draw, fill=white, diamond, aspect=1.9, align=center,
                   inner sep=0.5pt, text width=1.95cm,
                   execute at begin node={\hyphenpenalty=10000\relax}},
  term/.style   = {draw, fill=black!4, rounded corners=0.42cm, align=center,
                   text width=2.7cm, minimum height=0.84cm, inner sep=3pt,
                   execute at begin node={\hyphenpenalty=10000\relax}},
  con/.style    = {draw, circle, fill=white, minimum size=0.6cm, inner sep=0pt},
  nota/.style   = {draw, dashed, fill=white, align=center, text width=4.6cm,
                   inner sep=3pt, font=\fontsize{7.2}{8.6}\selectfont},
  marco/.style  = {draw, rounded corners=5pt, black!55, thin},
  tmarco/.style = {font=\fontsize{8}{9.6}\selectfont\itshape, anchor=north west},
  fl/.style     = {->, >=latex, semithick},
  ln/.style     = {semithick},
  et/.style     = {font=\fontsize{6.8}{8}\selectfont, inner sep=1.2pt, fill=white},
  union/.style  = {circle, fill=black, inner sep=0pt, minimum size=3pt}
}


\begin{document}
\begin{tikzpicture}[base]
% ======================= parte 1: captura y detección =======================
\node[term] (A) at (3.0,0)     {Inicio de impresión\\(estado por red local)};
\node[dec]  (F) at (3.0,-1.95) {¿Impresión terminada?};
\node[term, text width=2.3cm] (G1) at (6.35,-1.95) {Ciclo de remoción: etapas 1 a 4};
\node[proc] (B) at (3.0,-3.75) {Cada 10 min: iluminación encendida (D14), captura por Wi-Fi};
\node[proc] (D) at (3.0,-5.2) {Consulta al modelo de visión (API de OpenAI): captura y previsualización reducidas a 800 por 600};
\node[dec, text width=1.75cm]  (E) at (3.0,-7.1) {¿El modelo informa falla? (umbral por fijar)};
\node[negra] (H) at (3.0,-9.05) {Segunda consulta, imagen completa, para confirmar la falla};
\node[propp, text width=2.3cm] (S) at (6.35,-9.05) {Sigue sin supervisión, aviso al volver la red (P)};
\node[dec]  (I) at (3.0,-10.75) {¿Falla confirmada?};
\node[con]  (X1) at (3.0,-12.35) {2};

\draw[fl] (A) -- (F);
\draw[fl] (F.east) -- node[et, above] {sí} (G1.west);
\draw[fl] (F) -- node[et, right=1pt] {no} (B);
\draw[fl] (B) -- (D);
\draw[fl] (D) -- (E);
\draw[fl] (E) -- node[et, right=1pt] {sí} (H);
\draw[fl] (H) -- (I);
\draw[fl] (H.east) -- node[et, above] {sin red} (S.west);
\draw[fl] (I) -- node[et, right=1pt] {sí} (X1);
% retornos al inicio del lazo
\draw[ln] (E.west) -- node[et, above, pos=0.5] {no} (0.25,-7.1);
\draw[ln] (I.west) -- node[et, above, pos=0.5] {no} (0.25,-10.75) -- (0.25,-0.95);
\node[union] at (0.25,-7.1) {};
\draw[fl] (0.25,-0.95) -- (3.0,-0.95);
\draw[fl] (S.east) -- (7.85,-9.05) -- (7.85,-0.95) -- (3.0,-0.95);
\node[union] at (3.0,-0.95) {};

\draw[marco] (-0.1,1.0) rectangle (8.1,-12.95);
\node[tmarco] at (0.0,0.95) {Parte 1: captura y detección};

% ======================= parte 2: respuesta a la falla =======================
\node[con]  (X2) at (12.2,0)    {2};
\node[proc] (J)  at (12.2,-1.3) {Pausa de la impresora por red local (P)};
\node[proc] (K)  at (12.2,-2.75) {Aviso por Telegram: foto, descripción y botones};
\node[dec]  (L)  at (12.2,-4.4) {Decisión del usuario};
\node[proc, text width=2.35cm] (M) at (9.75,-6.35) {Reanudar la impresión (confirmación explícita)};
\node[term, text width=2.35cm] (R1) at (9.75,-8.05) {Vuelve a la captura (parte 1)};
\node[proc, text width=2.35cm] (N) at (12.4,-6.35) {Cancelar el trabajo: remoción de pieza fallida};
\node[term, text width=2.35cm] (G2) at (12.4,-8.05) {Ciclo de remoción: etapas 1 a 4};
\node[propp, text width=2.2cm] (O) at (15.0,-6.35) {Mantener la pausa y recordar cada 15 min (P)};

\draw[fl] (X2) -- (J);
\draw[fl] (J) -- (K);
\draw[fl] (K) -- (L);
\draw[fl] (L.west) -- node[et, above, pos=0.4] {reanudar} (9.75,-4.4) -- (M.north);
\draw[fl] (M) -- (R1);
\draw[fl] (L.south) -- node[et, right=1pt, pos=0.45] {cancelar} (12.2,-5.55) -- (12.4,-5.55) -- (N.north);
\draw[fl] (N) -- (G2);
\draw[fl] (L.east) -- node[et, above, pos=0.55] {sin respuesta} (15.0,-4.4) -- (O.north);
\draw[fl] (O.east) -- (16.4,-6.35) -- (16.4,-3.5) -- (12.2,-3.5);
\node[union] at (12.2,-3.5) {};

\draw[marco] (8.35,1.0) rectangle (16.65,-8.8);
\node[tmarco] at (8.45,0.95) {Parte 2: respuesta a la falla};
\end{tikzpicture}
\end{document}
